\section{Feature spaces and kernel machines}
\label{sec:kernels}

\subsection{Maximum-margin hyperplanes}

A linear classifier $h(x)=\mathrm{sign}(w\cdot x+b)$ on a linearly separable sample admits many separating hyperplanes. The \emph{margin} of $(w,b)$ is the distance to the nearest training point,
\begin{equation}
  \delta^\star=\min_{\mu\in[P]}\frac{\lvert w\cdot x_\mu+b\rvert}{\|w\|}.
  \label{eq:margin}
\end{equation}
The scale $(w,b)\sim(\alpha w,\alpha b)$ is redundant. Canonicalize by $\min_\mu\lvert w\cdot x_\mu+b\rvert=1$, so that $\delta^\star=1/\|w\|$. Maximizing the margin is then
\begin{equation}
  \min_{w,b}\ \tfrac12\|w\|^2
  \qquad\text{subject to}\qquad
  y_\mu(w\cdot x_\mu+b)\ge 1
  \quad\forall\mu.
  \label{eq:hard-margin}
\end{equation}

The Lagrangian is
\begin{equation}
  L(w,b,\alpha)
  =\tfrac12\|w\|^2
  -\sum_{\mu=1}^{P}\alpha_\mu\bigl(y_\mu(w\cdot x_\mu+b)-1\bigr),
  \qquad \alpha_\mu\ge 0.
  \label{eq:lagrange}
\end{equation}
Stationarity in $(w,b)$ gives the KKT conditions
\begin{equation}
  w=\sum_{\mu=1}^{P}\alpha_\mu y_\mu x_\mu,
  \qquad
  \sum_{\mu=1}^{P}\alpha_\mu y_\mu=0,
  \qquad
  \alpha_\mu\bigl(y_\mu(w\cdot x_\mu+b)-1\bigr)=0.
  \label{eq:kkt}
\end{equation}
Complementary slackness forces $\alpha_\mu=0$ unless the constraint is tight. Those points are the \emph{support vectors}; they alone determine $w$. Substituting~\eqref{eq:kkt} back into $L$ produces the dual
\begin{equation}
  \max_{\alpha\ge 0}\
  \sum_{\mu}\alpha_\mu
  -\tfrac12\sum_{\mu\nu}\alpha_\mu\alpha_\nu y_\mu y_\nu\, x_\mu\cdot x_\nu
  \qquad\text{s.t.}\qquad
  \sum_{\mu}\alpha_\mu y_\mu=0.
  \label{eq:svm-dual}
\end{equation}
The data enter only through the Gram matrix $G_{\mu\nu}=x_\mu\cdot x_\nu$. The classifier itself is
\begin{equation}
  h(x)=\mathrm{sign}\!\left(\sum_{\mu}\alpha_\mu y_\mu\, x\cdot x_\mu+b\right).
  \label{eq:svm-pred}
\end{equation}

\begin{figure}[ht]
\centering
\includegraphics[width=0.62\textwidth]{fig06_margin.pdf}
\caption{Hard-margin SVM in $\R^2$. Solid line: $w\cdot x+b=0$. Dashed: $w\cdot x+b=\pm 1$. Circled points are the support vectors.}
\label{fig:margin}
\end{figure}

\subsection{The kernel trick}

If $\{(x_\mu,y_\mu)\}$ is not linearly separable in $\R^{D}$, lift through a feature map $\Phi:\R^{D}\to\R^{M}$, $M\ge D$, and run the same SVM on $\Phi(x)$. The dual~\eqref{eq:svm-dual} and the predictor~\eqref{eq:svm-pred} depend on $\Phi$ only through
\begin{equation}
  K(x,x')\equiv\Phi(x)\cdot\Phi(x').
  \label{eq:kernel-def}
\end{equation}
One never instantiates $\Phi$. It is enough that $K$ be a symmetric positive-definite kernel (Mercer's theorem): there exist $\lambda_k\ge 0$ and orthonormal $\psi_k$ on $L^2(P_x)$ with
\begin{equation}
  K(x,x')=\sum_{k}\lambda_k\,\psi_k(x)\psi_k(x').
  \label{eq:mercer-kernel}
\end{equation}
The associated feature map is $\Phi(x)_k=\sqrt{\lambda_k}\,\psi_k(x)$, possibly infinite-dimensional.

\begin{theorem}[Representer]
Any minimizer of $\hat R_{\data}(h)+\lambda\|h\|_{\mathcal{H}_K}^2$ over the RKHS of $K$ is of the form
$h(\,\cdot\,)=\sum_{\mu=1}^{P}c_\mu K(x_\mu,\,\cdot\,)$.
\end{theorem}
The coefficients $c\in\R^{P}$ are the only degrees of freedom. Kernel ridge with the square loss~\eqref{eq:sq-loss} plus $\tfrac\lambda2\|w\|^2$ is derived in~\S\ref{sec:gen}:
\begin{equation}
  f(x)=k(x)^\top(\hat K+\lambda P I)^{-1} y,
  \qquad
  \hat K_{\mu\nu}=K(x_\mu,x_\nu),\quad
  k(x)_\mu=K(x,x_\mu).
  \label{eq:ridge}
\end{equation}

\subsection{A neural kernel}

A random one-hidden-layer net
\begin{equation}
  f(x)=N^{-1/2}\sum_{i=1}^{N} v_i\,\ph(w_i\cdot x),
  \qquad w_i\sim\mathcal{N}(0,D^{-1}I),\ \ v_i\sim\mathcal{N}(0,1),
  \label{eq:shallow}
\end{equation}
defines, at initialization, the kernel
\begin{equation}
  K_{\mathrm{NNGP}}(x,x')
  =\E_{w,v}\bigl[f(x)f(x')\bigr]
  =\E_{w}\bigl[\ph(w\cdot x)\ph(w\cdot x')\bigr].
  \label{eq:nngp-shallow}
\end{equation}
For $\ph(z)=\max(z,0)$ this is the Cho--Saul kernel. Derive it as follows. Let $w\sim\mathcal{N}(0,I)$, set $u=\hat x\cdot w$ and $u'=\hat x'\cdot w$ with $\hat x=x/\|x\|$, and write $\vartheta=\arccos(\hat x\cdot\hat x')$. Then $(u,u')$ is a standard bivariate Gaussian with correlation $\cos\vartheta$, equivalently the two coordinates of a 2d standard Gaussian $g$ in the plane of $(\hat x,\hat x')$:
\begin{equation}
  u=\rho\cos\phi,\qquad u'=\rho\cos(\phi-\vartheta),\qquad
  g=(\rho\cos\phi,\rho\sin\phi).
\end{equation}
ReLU is supported on $u>0$, $u'>0$, a wedge of opening $\pi-\vartheta$, namely $\phi\in(\vartheta-\pi/2,\pi/2)$. The Gaussian measure on the plane is $e^{-\rho^2/2}\,\rho\,\dd\rho\,\dd\phi/(2\pi)$, so
\begin{equation}
  \E[u_+u'_+]
  =\int_{\vartheta-\pi/2}^{\pi/2}\!\dd\phi
  \int_0^\infty\!\dd\rho\
  \frac{\rho^3 e^{-\rho^2/2}}{2\pi}\,\cos\phi\cos(\phi-\vartheta).
\end{equation}
The radial integral is $\int_0^\infty\rho^3 e^{-\rho^2/2}\,\dd\rho=2$. For the angle, $\cos\phi\cos(\phi-\vartheta)=\tfrac12\bigl(\cos(2\phi-\vartheta)+\cos\vartheta\bigr)$, so
\begin{equation}
  \int_{\vartheta-\pi/2}^{\pi/2}\cos\phi\cos(\phi-\vartheta)\,\dd\phi
  =\Bigl[\tfrac14\sin(2\phi-\vartheta)+\tfrac\phi2\cos\vartheta\Bigr]_{\vartheta-\pi/2}^{\pi/2}
  =\tfrac12\bigl(\sin\vartheta+(\pi-\vartheta)\cos\vartheta\bigr).
\end{equation}
Hence $\E[u_+u'_+]=(2\pi)^{-1}\bigl(\sin\vartheta+(\pi-\vartheta)\cos\vartheta\bigr)$. Scaling back $w\sim\mathcal{N}(0,D^{-1}I)$ multiplies by $\|x\|\|x'\|/D$:
\begin{equation}
  \E\bigl[\ph(w\cdot x)\ph(w\cdot x')\bigr]
  =\frac{\|x\|\|x'\|}{2\pi D}\bigl(\sin\vartheta+(\pi-\vartheta)\cos\vartheta\bigr).
  \label{eq:cho-saul}
\end{equation}
Deeper random nets compose this expectation layer by layer: that is the NNGP recursion below.

\begin{figure}[ht]
\centering
\includegraphics[width=0.72\textwidth]{fig07_chosaul.pdf}
\caption{Unit-variance Cho--Saul factor $\E[u_+u'_+]=(2\pi)^{-1}(\sin\vartheta+(\pi-\vartheta)\cos\vartheta)$ as a function of the angle between $x$ and $x'$.}
\label{fig:chosaul}
\end{figure}

\begin{exercise}
From~\eqref{eq:kkt}, show that if $x_\mu$ is not a support vector then $\alpha_\mu=0$. Conclude that $w$ is a linear combination of the support vectors only.
\end{exercise}
\begin{solution}
Complementary slackness in~\eqref{eq:kkt} is $\alpha_\mu\bigl(y_\mu(w\cdot x_\mu+b)-1\bigr)=0$. If $x_\mu$ is not a support vector the inequality is slack, $y_\mu(w\cdot x_\mu+b)>1$, so the second factor is nonzero and therefore $\alpha_\mu=0$. The stationarity condition $w=\sum_{\nu}\alpha_\nu y_\nu x_\nu$ then retains only those $\nu$ with $\alpha_\nu\neq 0$, i.e.\ the support vectors.
\ansbox{$x_\mu$ not on the margin $\Rightarrow$ $\alpha_\mu=0$, hence $w=\sum_{\mathrm{SV}}\alpha_\nu y_\nu x_\nu$.}
\end{solution}
